\Page{3}{\textenglish{ABSTRACT}}{אופנה פוגשת בינה מלאכותית}{36.5}{43}

{\fontsize{14.2}{20.4}\selectfont\bfseries
העין נמשכת לצבע, לגזרה או לתנועה של בד. משם קל להתאהב במראה שלם. הדרך שבה המראה הגיע אלינו עוברת במקומות שפועלים לפי היגיון אחר: שולחן פיתוח, מערכת תחזית, סט צילומים ומסך קניות. עכשיו נכנסים אליהם גם כלים של בינה מלאכותית.\par}

{\fontsize{13.1}{21.2}\selectfont
אופנה מזוהה עם דמיון, רגש וביטוי אישי. טכנולוגיה מחפשת חוקיות שאפשר להסביר, למדוד וליישם שוב. המפגש ביניהן מעניין במיוחד כשאחת מהן שואלת שאלה שהאחרת יכולה לעזור לברר. העין מזהה שינוי בצללית, ומערכת יכולה לבדוק אם הוא חוזר בתמונות ממקומות שונים. תוכנה יכולה להציע וריאציה נוספת, ומעצבת יכולה לשפוט אם היא מוסיפה משהו לשפה של הקולקציה.

ארבע הכתבות בגיליון בוחנות רגעים שונים שבהם החלטה יצירתית פוגשת חישוב. בכל רגע עומד למבחן דבר אחר: אם רעיון יכול להפוך לבגד, אם התעניינות יכולה להפוך לביקוש, אם דימוי יכול לספר סיפור אמין, ואם החלטה מוקדמת יכולה לשפר את מה שקורה לבגד בהמשך.

\textbf{עיצוב ופיתוח בגדים.} בין סקיצה לדוגמה ראשונה יש שורה של החלטות על בד, מבנה ומידות. בסביבת עבודה דיגיטלית אפשר לבדוק כמה כיוונים כשהבחירה עדיין פתוחה, ולנסח מוקדם יותר מה מבקשים מהיצרן. נראה איך נוצר תיק מוצר, מה אפשר ללמוד מהדמיה על גוף, ומה מתברר רק כשהבגד נתפר ונמדד. ככל שקל יותר לייצר הדמיות, עולה ערכה של השאלה אילו מהן מצדיקות עוד סיבוב של פיתוח.

\textbf{חיזוי טרנדים.} חזאיות סורקות תצוגות, חנויות ותרבות רחוב. תמונות ונתוני מסחר מאפשרים להן לעקוב גם אחרי קצב השינוי ולשאול איפה הוא מתרחש. הכתבה תבחן איך בודקים תחזית, ואיך קניינית שוקלת אם גוון שמושך את העין מצדיק הזמנה לעונה הבאה. לפריט חדש אפשר להקצות כמות קטנה, לנסות אותו בקהל מסוים או לשמור את הרעיון לעיתוי מתאים.

\textbf{קמפיינים ודמויות וירטואליות.} יום צילום תלוי בליהוק, בסטיילינג, בתאורה ובתיאום בין אנשים. כשחלק מההפקה נעשה באמצעות בינה מלאכותית, העבודה מתרחבת להכנת חומרי ייחוס, לבחירה בין גרסאות ולבדיקת התוצאה. נבחין בין ייצוג דיגיטלי של דוגמנית קיימת לבין דמות שנועדה להופיע שוב ושוב ולהחזיק חשבון משלה. לכל סוג הפקה יש שאלות משלו על שליטה בדימוי, על קרדיט ועל מה שהקהל מבין כשהוא מביט בתמונה. נשווה בין סוגי העבודה שנדרשים כדי להגיע לתוצר מוגדר, ונבדוק מה אפשר ללמוד מן העלויות שפורסמו והנתונים שנותרו בידיהם של המותגים.

\textbf{קנייה, ייצור ועתיד העבודה.} הדמיה על תמונת הלקוחה יכולה לעזור לה לבחור מה תרצה למדוד. בצד הייצור, מערכות תכנון יכולות להשפיע על הכמות ועל חלוקת המלאי בין חנויות. נבחן את ההבטחות לפי התוצאה המתאימה לכל החלטה, מבחירת מידה ועד למה שקורה למלאי שלא נמכר. כך גם מתברר איזה ידע עדיין צריך ליד המערכת: כיצד בד מתנהג, איך בודקים התאמה ומתי כדאי לחרוג מהמלצה. כדי להבין אם תכנון מדויק יותר יכול לצמצם בזבוז, צריך לעקוב גם אחרי המסלול שהבגד עובר מהחזרה לחנות, דרך מכירה חוזרת או שימוש מחדש, ועד למלאי שלבסוף מושמד.
\par}
